%%%PAGE 199 rot=0 legibility=fair summary=安培力F=kv分类讨论v-t图·标志重捕法与有机化学氢化物笔记·tanx导数图像·有空气阻力竖直上抛加速度·测电源电动势选电源
%%%BLOCK id=199-01 ch=12 sec=电磁感应·安培力F=kv分类讨论 kind=method alt=03 cont=none y=0.00-0.20
分类讨论\quad $kv$与$F_{\text{安}}$

\notefig[0.28]{figures/p199_a.png}

$F=kv$

\notefig[0.35]{figures/p199_b.png}

\notefig[0.6]{figures/p199_c.png}
% 图中标注：纵轴 $v$、横轴 $t$；三条曲线旁依次为 $k>\frac{B^2L^2}{R}$、$K=\frac{B^2L^2}{R}$（分母 $R$ 被涂改）、$k<\frac{B^2L^2}{R}$

$a=kv$\qquad $a=0$\qquad $v$\unclear{也为}$0$。
%%%END
%%%BLOCK id=199-02 ch=12 sec=电磁感应·i-t图像面积与电荷量 kind=note alt=none cont=none y=0.20-0.24
$i$-$t$图面积\ 若$\Delta\phi=0$\ 则面积$q$应为$0$。
%%%END
%%%BLOCK id=199-03 ch=99 sec=生物·标志重捕法与稀释涂布平板计数 kind=note alt=none cont=none y=0.27-0.42
% 与物理无关的生物笔记；后三行的“$\rho$小”被同一个椭圆圈起，首行“$\rho$大”下有一道横线
\begin{itemize}
\item 标志重捕法\ 标志物脱落\quad $\rho$大
\item 计数酵母菌时未稀释\quad $\rho$小
\item 在稀疏区取样\quad $\rho$小
\item 用稀释涂布平板计数大肠杆菌\quad $\rho$小
\end{itemize}
%%%END
%%%BLOCK id=199-04 ch=99 sec=化学·有机反应(内成环加成) kind=note alt=none cont=none y=0.45-0.50
% 与物理无关的化学笔记；红笔下划线共三处（“内成环”、“也”、“加成反应”，文字本身为黑色）
分子\red{\uline{内成环}}\red{\uline{也}}叫\red{\uline{加成反应}}，（啥也没加）。
%%%END
%%%BLOCK id=199-05 ch=99 sec=化学·氢化物与硅硫硼化合物 kind=note alt=none cont=none y=0.51-0.69
% 与物理无关的化学笔记。红笔：首句“氢化”下、方框内“最简单氢化物”下、两个圆圈内文字下各有一道红色下划线；方框内压着一个五角星
\red{\uline{氢化}}物注意\ \fbox{最简单氢化物 ★}\ \ \fbox{CH$_4$}\ 与\ \fbox{\unclear{沥青}}\ 都是氢化物

\notefig[0.6]{figures/p199_d.png}
% 结构式转录：(CH$_3$)$_3$ *Si$-$S$-$B，B 上分两支：上支 $-$S$-$Si(CH$_3$)$_3$；下支 S \struck{CH$_3$} 另起 Si$-$(CH$_3$)$_3$
%%%END
%%%BLOCK id=199-06 ch=90 sec=微积分·常用导数(tan x) kind=formula alt=none cont=none y=0.69-0.86
$(\tan x)'=\dfrac{1}{\cos^2 x}$

\notefig[0.55]{figures/p199_e.png}
% 图为 $y=\tan x$ 的导函数 $\frac{1}{\cos^2x}$ 的草图（竖直虚线为渐近线），箭头自公式指向中间一支；右侧另有一小段红笔画的 V 形曲线
%%%END
%%%BLOCK id=199-07 ch=03 sec=有空气阻力的竖直上抛 kind=formula alt=01 cont=none y=0.86-0.95
有空气阻力

\notefig[0.1]{figures/p199_f.png}
% 图为两个带小圆圈（物体）的竖直箭头：上方箭头向上，下方箭头向下

\[
a_1=\frac{f+mg}{m}
\]
\[
a_2=\frac{mg-f}{m}\ \,{}^{a\text{不同}}
\]
%%%END
%%%BLOCK id=199-08 ch=10 sec=测电源电动势和内阻·电源选择 kind=note exp=1 alt=none cont=none y=0.95-1.00
测电源电动势、内阻实验中\qquad $E=I_{\max}R_{\text{滑}\max}$

选$E$略大于$I_{\max}R_{\text{滑}\max}$的电源
% 后一行为写在上一行右上方的小字批注，一个弯箭头从该批注指向上一行的“$E$”；“$I_{\max}R_{\text{滑}\max}$”下有红笔下划线
%%%END
%%%PAGEEND 199
